% -----------------------------------------------------------
\section{Round}
% -----------------------------------------------------------
	\label{ROUND}
	
% % ----------------------
% \subsection{See also}
% % ----------------------

% ----------------------
\subsection{1828 American Dictionary of the English Language} % *1828
% ----------------------
	\label{ROUND-1828}
	
	% -----
	\subsubsection{Noun}
	% -----

		\begin{compactitem}
			\item[1] A circle; a circular thing, or a circle in motion.
			\item[2] Action or performance in a circle, or passing through a series of hands or things and coming to the point of beginning; or the time of such action.
			\item[3] Rotation in office; succession in vicissitude.
			\item[4] A \hyperref[RUNDLE]{rundle}; the step of a ladder.
				\begin{compactitem}
					\item Compare OED \#III.13.a.
				\end{compactitem}
			\item[5] A walk performed by a guard or an officer.
			\item[6] A dance; a song; a roundelay, or a species of fugue.
			\item[7] A general discharge of fire-arms by a body of troops, in which each soldier fires once.
		\end{compactitem}

% ----------------------
\subsection{Oxford English Dictionary} % *OED
% ----------------------
	\label{ROUND-OED}
	
	% -----
	\subsubsection{Noun}
	% -----

		\begin{compactitem}
			\item[I] Something circular in form, and related senses; a circle, a ring.
			\item[I.1.a] A circular band, a ring; a ring-shaped object; (in later use) spec. a representation or drawing of a ring-shaped object on a flat surface, esp. one used to enclose another image, an inscription, etc. Formerly also: †a coil (obsolete). Also figurative.
			\item[I.1.b] A link in a chain. Obsolete.
			\item[I.1.c] A single turn of thread, yarn, etc., wound on to a reel.
			\item[I.2.a] A group of people sitting (standing, dancing, etc.) in a circle. Frequently in in a round.
			\item[I.2.b] A circular group of things; a number of things set or arranged in a ring.
			\item[I.2.c] In extended use: a group or set of things which are related or typically found together, esp. a circle of friends or acquaintances.
			\item[I.3] The circumference, bounds, or extent of a circular object. Also figurative. Now chiefly literary.
			\item[I.4.a] A circular object, a disc; spec. a representation or drawing of a circle on a flat surface, e.g. a heraldic roundel. 
			\item[I.4.b] A circular archery target marked with a bullseye at the centre of a number of concentric rings, and usually mounted on a butt (butt n.7 2a). Chiefly in at rounds. Frequently contrasted with rover n.2 1a. Obsolete (historical in later use).
			\item[I.4.c] A flat circular piece of something. Frequently in Cookery.
			\item[I.4.d] A large round piece of beef; spec. a thick disc cut from the haunch as a joint. Frequently in round of beef.
			\item[I.4.e] A slice of bread, originally as cut from a round loaf; a slice of toast.
			\item[I.5.a] A man-made structure with a circular form; a curved or rounded part of such a structure; (in early use) esp. a circular tower or turret forming part of a larger structure. Chiefly Scottish in early use.
			\item[I.5.b] A natural structure of circular form; (now) esp. a ring or circular clump or flowers, trees, etc. Also: a circular or rounded part of a natural structure.
			\item[I.5.c] A curve or bend in a river, coastline, etc. Now rare.
			\item[I.6] A statue or piece of sculpture made in the round (see Phrases 2a(a)). Also: a fragment of a statue. Obsolete.
			\item[I.7] Brewing. Originally (more fully cleansing round): a large cask in which beer is purified and the fermentation process finished (now chiefly historical). Later (usually in fermenting round): a large round vessel in which beer is fermented.
			\item[I.8] Surveying. A complete circle (of measured angles). Now rare.
			\item[II] A spherical object.
			\item[II.9] Something of spherical form; a sphere, an orb; spec. a planet. Formerly also: a globular lump (obsolete). Somewhat rare except as in sense 10a.
			\item[II.10] spec. Chiefly poetic.
			\item[II.10.a] With this (also the). The earth. Frequently with modifying word, esp. in this earthly round, this terrestrial round.
			\item[II.10.b] The vault of the sky; the firmament; = heaven n. 1a. Frequently more fully round of heaven, celestial round, etc.
			\item[II.11] = round shot n.   Now chiefly historical.
			\item[II.12] A (more or less) spherical variety of potato.
			\item[III] A cylindrical object, and related senses.
			\item[III.13] A cylindrical part of an implement or other manufactured object.
			\item[III.13.a] A rung of a ladder.
			\item[III.13.b] figurative or in figurative context. Cf. rung n. 2b. Now rare.
			\item[III.13.c] A rod or bar joining and supporting the legs of a chair (or occasionally the stilts of a plough); a stretcher. rare before 19th cent. (chiefly U.S. in later use).
			\item[III.13.d] Printing. = rounce n.1 1. Obsolete. rare.
			\item[III.13.e] Each of the cylindrical pins of a trundle or lantern pinion.
			\item[III.14] A long length of fabric rolled or wound up into cylindrical form; this as a fixed measurement, esp. for sale. Cf. roll n.1 5a. Obsolete.
			\item[III.15] An object with a convex surface.
			\item[III.15.a] Architecture. A convex moulding of which the section is (more or less) a semicircle, as an astragal, a torus, etc.
			\item[III.15.b] Woodworking. A plane with a convex iron designed to cut an arc of sixty degrees, typically used in conjunction with a hollow.
			\item[III.16] A cylindrical piece of something.
			\item[III.16.a] A roughly cylindrical piece of edible root, tuber, or stalk, esp. of dried rhubarb rhizome.
			\item[III.16.b] An iron or steel bar of circular section.
			\item[IV] Something which moves in a circular manner or occurs cyclically.
			\item[IV.17.a] A dance in which the dancers move in a circle, typically while remaining in contact with each other, e.g. by holding hands. Also figurative and in extended use.
			\item[IV.17.b] The music for such a dance. rare.
			\item[IV.18] A swinging stroke or blow with a sword. Obsolete.
			\item[IV.19] Music. Originally: a short, simple song, esp. one in which two or more people sing in turn. Later: spec. a song for three or more unaccompanied voices or parts, each singing the same theme but starting one after another, at the same pitch or in octaves; a simple canon.
			\item[IV.20.a] A regularly recurring sequence or continuous cycle of activities, events, duties, etc. Frequently with of.
			\item[IV.20.b] The constant passage and recurrence of days, years, or other cyclical measures of time; the passage of time in cyclical periods. Also: a particular cycle of days, hours, etc. Usually with of.
			\item[IV.20.c] Theosophy. Each of seven subdivisions of human spiritual life, represented as a cycle of spiritual development.
			\item[IV.21] An act of ringing each bell of a set or peal in sequence; (now) spec. a sequence in which bells are rung moving down the scale in order, typically preceding the ringing of changes (change n. 9a). Also occasionally in extended use.
			\item[V] A circular route; a circuit, a tour, and related senses. Frequently in combination with a basic verb of action, as go, make, take, etc.
			\item[V.22.a] A circular or revolving motion; movement in a circle; esp. a recurring circular course. Also in plural and figurative.
			\item[V.22.b] An indirect or circuitous path, route, or journey; a journey which contains many detours and ultimately comes full circle. Now somewhat rare.
			\item[V.22.c] A circuit of a racecourse; a lap.
			\item[V.23] \textit{Military}.
			\item[V.23.a] A circuit of a garrison or camp, the ramparts of a fortress, etc., made by a patrol, esp. during the night, to ensure that the sentinels are vigilant; (also) a circuit of the streets of a town made by a watch to preserve good order. Chiefly in to go (also †make, pace, walk, etc.) the rounds (also round). Also figurative and in extended use. Now chiefly in plural and historical.
			\item[V.23.b] A watch or patrol that is responsible for making a regular inspection of sentries, patrolling a town, etc. Now rare.
			\item[V.23.c] Nautical. In plural. Inspection.
			\item[V.24] In singular or plural. Frequently with possessive adjective, as to go (also walk, march, etc.) one's round (also rounds).
			\item[V.24.a] A route, course, or circuit habitually used or followed by a person in the course of his or her job. 
			\item[V.24.b] \textit{spec}.
			\item[V.24.b.(a)] Chiefly British. A journey along a fixed route for the purpose of delivering goods; a regular route followed in this way; a delivery job involving journeys along such a route.
			\item[V.24.b.(b)] A regular visit by a doctor or nurse in a hospital to each of the patients under his or her care.
			\item[V.25] The circuit of a place, building, etc. Now only with of. Also figurative.
			\item[V.26.a] A recreational walk or drive round a place or a number of places; a tour. Usually in to take (also go, make) a round. Also figurative. Now chiefly South Asian.
			\item[V.26.b] An act of visiting a number of people or places in turn. Frequently in plural.
			\item[V.26.c] A series of visits, appearances, etc., undertaken in a single trip or tour, or within a certain period of time.
			\item[V.27] In plural. With the. The procedure established under the Poor Law by which unemployed agricultural labourers were given temporary work by a number of farmers in turn. Chiefly in on the rounds, (also) to go the rounds. Cf. roundsman n. 1. Now historical.
			\item[V.28] Australian and New Zealand. Frequently with modifying word. Originally: the succession of routine enquiries made by a journalist covering news stories in a particular subject or area (frequently in plural); esp. in police round. Hence also: the subject or area assigned for a journalist to cover. 
			\item[VI] A single amount or quantity of something.
			\item[VI.29.a] A set of drinks bought for all the members of a group, in early use esp. for a toast, now typically as one of a sequence bought by different members of the group. Formerly frequently with of and type of drink specified.
			\item[VI.29.b] A set of portions of food bought or provided for all the members of a group. Frequently with of and type of food specified.
			\item[VI.29.c] Chiefly British. The quantity of sandwiches (typically two) made from two slices of bread.
			\item[VI.30] Frequently in \textit{plural}.
			\item[VI.30.a] The amount of ammunition needed to fire one shot from a firearm or piece of artillery. Also occasionally: †a single arrow or bolt for a bow (obsolete).
			\item[VI.30.b] A single volley of fire by artillery or a number of firearms; a shot from a single firearm or piece of artillery. Also: †a single arrow shot from a bow (obsolete). 
			\item[VI.31] Coal Mining. A quantity of coal delivered to the surface by a putter. Obsolete.
			\item[VI.32] A single, distinct outburst of applause, cheers, etc.
			\item[VI.33] A unit of trade used by European traders in West Africa, consisting of a selection of goods. Frequently in round trade n. at Compounds 2. Now historical.
			\item[VII] A period or bout of play at a game or sport, and related senses.
			\item[VII.34.a] A game of cards. Chiefly with of.
			\item[VII.34.b] Golf. An act of playing all the holes on a golf course (or occasionally all those nominated for play in a particular tournament) once. Also (with of): the score made in doing this.
			\item[VII.34.c] gen. A bout of play at a game, sport, or contest; (in early use) spec. a bout of fisticuffs. Also in extended use.
			\item[VII.35] Cards. A single turn of play by all the players, esp. the playing of a single card in turn by each of the players; a trick.
			\item[VII.36.a] Originally Boxing. Each of the periods into which a boxing match (or later a match in other combat sports, as wrestling, karate, etc.) is divided. Also in extended use.
			\item[VII.36.b] figurative. With reference to the exchanges and contributions constituting a debate, argument, controversy, etc. Frequently in round two, typically suggesting the initiation of a new phase of the dispute.
			\item[VII.36.c] Each of a number of subdivisions or segments of a quiz, game show, or similar competition.
			\item[VII.37] Originally and chiefly Sport. Each of a succession of stages in a competition, each resulting in the elimination of a number of competitors or candidates.
			\item[VII.38] Archery. A fixed number of arrows shot from a fixed distance. Frequently with distinguishing word, as National round, Vegas round, York round, etc.
			\item[VII.39] A session of meetings for discussion or negotiation, esp. one of a sequence of sessions in an ongoing process, typically characterized by a development between one session and another.
		\end{compactitem}

% % ----------------------
% \subsection{Scriptural Usage} % *SCRIPTURES
% % ----------------------
	% \label{ROUND-SCRIPTURES}

	% % -----
	% \subsubsection{Old Testament} % *OT
	% % -----
		% \label{ROUND-OT}
	
		% % --
		% \paragraph{KJV}
		% % --
			% \label{ROUND-OT-KJV}
	
			% \begin{tabular}{ll}
				% Total occurrences in the OT & 
			% \end{tabular}
			
			% \begin{compactdesc}
				% \item[]
			% \end{compactdesc}
			
		% % --
		% \paragraph{Hebrew Bible}
		% % --
			% \label{ROUND-HEBREW}
		
			% %
			% \subparagraph{\texorpdfstring{\he{sample word}}{}}
			% %
				% \label{PAGA} % whatever is in step bible without periods
			
				% \begin{compactdesc}
					% \item[HALOT , PDF ] \textbf{}
						% \begin{compactitem}
							% \item[1]
						% \end{compactitem}
					% \item[BDB , PDF ] \textbf{}
						% \begin{compactitem}
							% \item[1]
						% \end{compactitem}
					% \item[DCH PDF ] \textbf{}
						% \begin{compactitem}
							% \item[1]
						% \end{compactitem}
				% \end{compactdesc}
				
				% \begin{compactdesc}
					% \item[Some OT uses] \textbf{}
						% \begin{compactdesc}
							% \item[]
						% \end{compactdesc}
				% \end{compactdesc}
			
		% % --
		% \paragraph{Septuagint}
		% % --
			% \label{ROUND-LXX}
		
			% %
			% \subparagraph{\texorpdfstring{\gr{sample word}}{}}
			% %
		
	% % -----
	% \subsubsection{New Testament} % *NT
	% % -----
		% \label{ROUND-NT}
	
		% % --
		% \paragraph{KJV}
		% % --
			% \label{ROUND-NT-KJV}
			
	
			% \begin{tabular}{ll}
				% Total occurrences in the NT & 
			% \end{tabular}
			
			% \begin{compactdesc}
				% \item[]
			% \end{compactdesc}
			
		% % --
		% \paragraph{Greek New Testament} % *GREEK
		% % --
			% \label{ROUND-GREEK}
		
			% %
			% \subparagraph{\texorpdfstring{\gr{sample word}}{}}
			% %
				% \label{KATALLAGE}
			
				% \begin{compactdesc}
					% \item[BDAG , PDF ] \textbf{}
						% \begin{compactitem}
							% \item 
						% \end{compactitem}
					% \item[LSJ , PDF ] \textbf{}
						% \begin{compactitem}
							% \item 
						% \end{compactitem}
				% \end{compactdesc}
				
				% \begin{compactdesc}
					% \item[Some NT uses] \textbf{}
						% \begin{compactdesc}
							% \item[]
						% \end{compactdesc}
				% \end{compactdesc}
		
	% % -----
	% \subsubsection{Book of Mormon} % *BOM
	% % -----
		% \label{ROUND-BOM}
	
		% \begin{tabular}{ll}
			% Total occurrences in the BOM & 
		% \end{tabular}
		
		% \begin{compactdesc}
			% \item[]
		% \end{compactdesc}
		
	% % -----
	% \subsubsection{Doctrine and Covenants} % *DC
	% % -----
		% \label{ROUND-DC}
	
		% \begin{tabular}{ll}
			% Total occurrences in the DC & 
		% \end{tabular}
		
		% \begin{compactdesc}
			% \item[]
		% \end{compactdesc}
		
	% % -----
	% \subsubsection{Pearl of Great Price} % *PGP
	% % -----
		% \label{ROUND-PGP}
	
		% \begin{tabular}{ll}
			% Total occurrences in the PGP & 
		% \end{tabular}
		
		% \begin{compactdesc}
			% \item[]
		% \end{compactdesc}
		
% ----------------------
\subsection{Key Phrases} % *KEYPHRASES *PHRASES
% ----------------------
	\label{ROUND-PHRASES}

	\begin{compactdesc}
	
		\item[one eternal round] \textbf{5} | \hyperref[1NEPHI-10-19]{1 Nephi 10:19}; \hyperref[ALMA-7-20]{Alma 7:20}; \hyperref[ALMA-37-12]{37:12}; \hyperref[DC3-2]{DC 3:2}; \hyperref[DC35-1]{DC 35:1}
			\begin{compactdesc}
				% \item[Bible anchor verses] \phantom{.}
					% \begin{compactdesc}
						% \item[Romans 5:5] \gr{}
					% \end{compactdesc}
				\item[Commentary] This is a very interesting phrase; notice that all five occurrences of it are describing the ``course'' of the Lord.
				
					Here is some more material for thinking about this phrase:
					
					Some definitions of ``course'' in the 1828 dictionary:
						\begin{compactitem}
							\item[1] In its general sense, a passing; a moving, or motion forward, in a direct or curving line; applicable to any body or substance, solid or fluid.
							\item[2] The direction of motion; line of advancing; point of compass, in which motion is directed.
							\item[4] A passing or process; the progress of any thing; as the course of an argument, or of a debate; a course of thought or reflexion.
							\item[5] Order of proceeding or of passing from an ancestor to an heir; as the course of descent in inheritance.
							\item[6] Order; turn; class; succession of one to another in office, or duty.
							\item[7] Stated and orderly method of proceeding; usual manner. 
							\item[8] Series of successive and methodical procedure; a train of acts, or applications.
							\item[9] A methodical series, applied to the arts or sciences; a systemized order of principles in arts or sciences, for illustration of instruction.
							\item[10] Manner of proceeding; way of life or conduct; deportment; series of actions.
							\item[11] Line of conduct; manner of proceeding.
						\end{compactitem}
					There are also various meanings of ``\hyperref[ROUND]{round}'' that give various senses of the phrase. Here are a few from the OED, along with an explanation of what the phrase might mean with that sense of ``\hyperref[ROUND]{round}'':
						\begin{compactitem}
							\item[I] The motion the Lord is taking forward, either literally or figuratively, repeats itself again and again, as in a circle.
							\item[III.13.a] (``The course the Lord is taking in this implementation of the plan of salvation is one step up the eternal ladder of progression.'') The progress the Lord makes is continually upward, one step at a time, as though he were ascending a ladder.
							\item[IV] The progress the Lord undergoes occurs cyclically, passing through the same or similar stages of development again and again (see also IV.20.a and IV.20.b).
						\end{compactitem}
					All three of these, but especially those in IV, have a natural Adam-God interpretation, where the cycle of resurrection, exaltation, fall, progression, resurrection is followed again and again. I think this is an extremely powerful and beautiful interpretation.
					
					From the scriptures themselves we can also gather more about possible meanings of this phrase:
						\begin{compactitem}
							\item[1 Nephi 10:19] God works in the same way with his children across time; in particular, the Holy Ghost reveals things to them regardless of what time period they live in.
							\item[Alma 7:20, Alma 37:12, DC 3:2] The Lord takes the straightest and most efficient route to complete anything. Also, he does not vary from the course to go in some aberrant direction.
							\item[DC 35:1] Similar to 1 Nephi 10:19, the Lord uses the smae methods of acting with his children throughout time.
						\end{compactitem}
						
					Another verse to consider here is \hyperref[LUKE-9-31]{Luke 9:31}, which uses the term \gr{ἔξοδος}, which can mean ``course'' but also ``death.'' So ``one eternal round'' on this meaning, if it were understood as the ``course,'' could be the Adam-God idea of repeating the cycle of dying, beng born, and being exalted again.
					
					This phrase is really rich with meaning. I'm sure that I have not exhausted the possibilities here.
			\end{compactdesc}
			
	\end{compactdesc}
	
% % ----------------------
% \subsection{Bible Dictionaries} % *DICTIONARIES
% % ----------------------
	% \label{ROUND-BD}

% % ----------------------
% \subsection{Restoration Usage} % *RESTORATION
% % ----------------------
	% \label{ROUND-RESTORATION}

	% % -----
	% \subsubsection{Joseph Smith} % *JS
	% % -----
		% \label{ROUND-JS}
	
		% \begin{compactdesc}
			% \item[{\hyperref[KEY]{Date/Source}}]
		% \end{compactdesc}
	
	% % -----
	% \subsubsection{Brigham Young} % *BY
	% % -----
		% \label{ROUND-BY}
	
		% \begin{compactdesc}
			% \item[{\hyperref[KEY]{Date/Source}}]
		% \end{compactdesc}
	
% % ----------------------
% \subsection{Commentary and Studies} % *COMMENTARY *STUDIES
% % ----------------------
	% \label{ROUND-COMMENTARY}

	% % -----
	% \subsubsection{Key Points}
	% % -----
		% \label{ROUND-KEYS}
	
		% \begin{compactitem}
			% \item
		% \end{compactitem}
		
	% % -----
	% \subsubsection{Sample Study}
	% % -----
		% \label{ROUND-study}